\subsection{Problem Setup}
\label{sec:problem_setup}

We study \emph{fixed-pool answer selection under a matched-budget inference
contract}.
The task is: given a set of candidate answers already produced by a collection
of inference methods under an equal per-method budget, select the best final
answer without issuing additional provider calls and without access to
ground-truth labels at decision time.

\paragraph{Problem instance.}
Each instance comprises a question $q$, a \emph{candidate pool}
$\mathcal{C}(q) = \{a_F, a_{L1}, a_{S1}, a_{T}\}$ of four final-answer
candidates generated under the matched-budget protocol
(Section~\ref{sec:experiments}), and failure-trace metadata $m_F(q)$ produced
as a byproduct of the frontier method's normal execution.
A ground-truth answer $y(q)$ exists but is withheld at runtime; it is used only
for offline evaluation.
The selection policy maps $(\mathcal{C}(q),\, m_F(q)) \to \hat{a}(q) \in
\mathcal{C}(q)$.

\paragraph{Candidate pool and budget accounting.}
The candidate pool is constructed once per evaluation run.
The four candidate-generating methods each produce one final-answer candidate per
example under a cap of $B = 6$ logical API calls: the frontier method and three
external baselines (L1, S1, and TALE).
Full candidate-pool construction therefore requires $4 \times B = 24$ logical
inference calls per example.
The answer-selection policy (FTA) is applied \emph{after} pool construction:
it selects among already-generated candidates and issues \emph{no} additional
inference calls.

\paragraph{Runtime-visible vs.\ offline-only information.}
All inputs to the selection policy are runtime-visible and gold-free: the
frontier answer, the three external answers, and failure-trace metadata derived
from the frontier method's logged execution.
Ground-truth labels, exact-match signals, and correctness indicators are
\emph{never} inputs to any policy decision; they are used only offline for
evaluation and failure analysis.
The oracle accuracy (92.00\% on Final-300) is a diagnostic upper bound computed
offline with gold labels; it is not a comparison baseline.

Evaluation results are reported on:
\begin{itemize}[nosep,leftmargin=1.5em]
    \item Final-300 (seed~71, $n=300$): the primary unbiased evaluation set,
    \item Aggregate-720 (seeds 41, 61, 71; $n=300+120+300$): a broader
          disjoint aggregate used as supporting robustness evidence within the
          same experimental regime, with zero example-id and question-hash
          overlap across sources,
    \item Seed~61 ($n=120$): a failure-enriched development split; absolute
          accuracy on this split is not representative of standard evaluation
          conditions and is not used to derive primary claims.
\end{itemize}

\subsection{Terminology}
\label{sec:terminology}

The following terms are used consistently throughout the paper.
Short paper-facing names are used for readability; the corresponding
implementation identifiers are preserved in Table~\ref{tab:name_mapping} for
reproducibility.

\begin{description}[leftmargin=2.2cm, labelwidth=2cm, labelsep=0.2cm, style=nextline]
  \item[Candidate pool] The set of final answers produced by the available
    methods (frontier and external baselines) for a given example.
    Candidate-pool construction is separate from the downstream selection
    policy.
  \item[Frontier method] The primary inference method whose answer serves as
    the default output before any policy gate fires.
    It produces failure-trace metadata as a side-effect of normal execution,
    including structural indicators of search breadth and internal answer
    agreement.
  \item[Failure-trace metadata] Runtime-visible fields logged during frontier
    execution that characterize the structural quality of the frontier's
    search output (e.g., whether the selected answer arises from a single weak
    branch or from an internally consistent exploration).
    These fields are gold-free by construction.
  \item[External baseline] Any of the three independently-sourced single-model
    methods used as reference signals: L1, S1, or TALE.
    External baselines provide gold-free consensus signals without access to
    gold labels.
  \item[L1] An external method that selects the answer with the highest support
    count across sampled model outputs, following the self-consistency answer
    aggregation principle \citep{wang2023selfconsistency}.
  \item[S1] An external method that applies sequential budget-forcing to elicit
    a final answer within the allocated budget \citep{muennighoff2025s1}.
  \item[TALE] An external method that uses prompt-level budget cues to guide
    answer selection \citep{han2025token}.
  \item[Logical API call] A single provider inference request counted against a
    candidate-producing method's budget.
  \item[Budget] The integer upper bound $B = 6$ on logical API calls per
    example for each candidate-producing method, applied identically to all
    compared methods (matched-budget condition).
  \item[Matched budget] All compared candidate-producing methods operate under
    the same per-method budget cap $B$, ensuring that accuracy differences are
    not confounded by unequal per-method computational limits.
  \item[Gold-free runtime feature] Any input to a policy decision rule derived
    solely from model outputs, run metadata, or structural properties of the
    candidate pool---never from gold labels, exact-match signals, or
    correctness indicators.
  \item[Oracle upper bound] The accuracy achievable by a hypothetical policy
    with offline access to gold labels, used only as a diagnostic ceiling.
    It is not a baseline method.
  \item[Provider-specific evidence] Results obtained from a particular
    provider--model pairing (here, Cohere command-r-plus-08-2024).
    Claims are explicitly scoped to this setting and do not generalize beyond
    it without additional replication.
\end{description}

These three external baselines are run under the same $B = 6$ logical-call cap
as the frontier method. L1 is a support-count, self-consistency-inspired answer
aggregation baseline; S1 adapts the budget-forcing inference idea from s1; and
TALE is a token-budget-aware prompting baseline. Together they serve as
independently sourced external reference signals for FTA.
